% Generated draft: requires mathematical and source review.
\documentclass[11pt]{article}
\usepackage[margin=25mm]{geometry}
\usepackage{fontspec}
\setmainfont{Latin Modern Roman}
\usepackage{amsmath,amssymb,mathtools,unicode-math}
\setmathfont{Latin Modern Math}
\usepackage{enumitem,needspace,xurl,hyperref}
\setlength{\parindent}{0pt}
\setlength{\parskip}{5pt}
\setlength{\emergencystretch}{3em}
\newcommand{\N}{\mathbb{N}}
\newcommand{\Z}{\mathbb{Z}}
\newcommand{\Q}{\mathbb{Q}}
\newcommand{\R}{\mathbb{R}}
\newcommand{\C}{\mathbb{C}}
\newcommand{\F}{\mathbb{F}}
\newcommand{\A}{\mathbb{A}}
\newcommand{\PP}{\mathbb{P}}
\newcommand{\E}{\mathbb{E}}
\newcommand{\eps}{\varepsilon}
\newcommand{\dd}{\,\mathrm{d}}
\begin{document}
\section*{UnsolvedMath ID: 30002517}
\textbf{Source problem code:} catalogue entry
\subsection{Definitions and mathematical statement}
Write \(\mathbf x\) for the initial cluster, \(Q\) for its quiver, and \(\mathcal A(\mathbf x,Q)\) for the associated cluster algebra under the intended seed and coefficient conventions. The supplied record does not specify the ground ring, coefficient system, frozen vertices, which coefficients are inverted, or the admissible class of quivers. These choices remain unresolved and are not fixed here.

For clarity, if \(A=\mathcal A(\mathbf x,Q)\) is a Krull domain, let \(K=\operatorname{Frac}(A)\), let \(X^{(1)}(A)\) denote its set of height-one prime ideals, and let \(v_{\mathfrak p}:K^\times\to\Z\) be the normalized discrete valuation associated with \(A_{\mathfrak p}\). Its divisor group and divisor class group are
\[
\operatorname{Div}(A)=\bigoplus_{\mathfrak p\in X^{(1)}(A)}\Z[\mathfrak p],\qquad
\operatorname{Cl}(A)=\operatorname{Div}(A)\big/\{\operatorname{div}(f):f\in K^\times\},
\]
where
\[
\operatorname{div}(f)=\sum_{\mathfrak p\in X^{(1)}(A)}v_{\mathfrak p}(f)[\mathfrak p].
\]
The Krull-domain hypothesis makes this standard divisor-theoretic interpretation available; it is not stated as a hypothesis of the supplied question.

\textbf{Problem.} Determine the divisor class group
\[
\operatorname{Cl}\bigl(\mathcal A(\mathbf x,Q)\bigr)
\]
of the cluster algebra associated with \(Q\). A fully quantified formulation requires the missing seed, coefficient, and quiver conventions. If cluster algebras outside the Krull-domain setting are intended, the meaning of their divisor class group must also be specified.

The supplied dated curation labels the problem partially solved. It reports explicit computability for acyclic seeds and finite generation and freeness of the divisor class group for every cluster algebra that is a Krull domain, while reporting no established uniform answer in the stated generality. These are supplied literature assessments, not independently verified conclusions of this draft.
\subsection{Short English statement}
Determine the divisor class group of a cluster algebra from its initial cluster and quiver. The source leaves the coefficient conventions and permitted quivers unspecified, and does not assume that the algebra is a Krull domain. The standard Krull-domain interpretation is given conditionally; the intended full generality requires clarification.
\subsection{Sources}
\begin{itemize}
\item Supplied local UnsolvedMath registry record, permanent ID \(30002517\), title \emph{Divisor Class Groups of Cluster Algebras}, source problem code \texttt{catalogue entry}, in the Oberwolfach Reports Open Problems collection. The record gives the proposed year as 2014 and identifies no individual proposer.
\item \emph{Representation Theory of Quivers and Finite Dimensional Algebras} (2014), Oberwolfach Reports, pp.\ 453--529, DOI \texttt{10.4171/owr/2014/08}. Supplied source and literature reference: \url{https://doi.org/10.4171/owr/2014/08}. The supplied background describes these reports as documenting mathematical workshops and their open-problem discussions.
\item \emph{Factoriality and class groups of cluster algebras}. Supplied literature reference: \url{https://arxiv.org/abs/1712.06512}.
\item Literature assessment checked at \texttt{2026-08-21}: the supplied curation labels the problem partially solved, reports explicit computability for acyclic seeds and finite generation and freeness for cluster algebras that are Krull domains, and states that a uniform answer in the stated generality is not established. This is supplied dated curation derived from the OWR record, not an independently verified literature assessment; its status and cited claims require independent verification.
\end{itemize}
\end{document}
